% Generated draft: requires mathematical and source review.
\documentclass[11pt]{article}
\usepackage[margin=25mm]{geometry}
\usepackage{fontspec}
\setmainfont{Latin Modern Roman}
\usepackage{amsmath,amssymb,mathtools,unicode-math}
\setmathfont{Latin Modern Math}
\usepackage{enumitem,needspace,xurl,hyperref}
\setlength{\parindent}{0pt}
\setlength{\parskip}{5pt}
\setlength{\emergencystretch}{3em}
\newcommand{\N}{\mathbb{N}}
\newcommand{\Z}{\mathbb{Z}}
\newcommand{\Q}{\mathbb{Q}}
\newcommand{\R}{\mathbb{R}}
\newcommand{\C}{\mathbb{C}}
\newcommand{\F}{\mathbb{F}}
\newcommand{\A}{\mathbb{A}}
\newcommand{\PP}{\mathbb{P}}
\newcommand{\E}{\mathbb{E}}
\newcommand{\eps}{\varepsilon}
\newcommand{\dd}{\,\mathrm{d}}
\begin{document}
\section*{UnsolvedMath ID: 11100054}
\textbf{Source problem code:} catalogue entry
\subsection{Definitions and mathematical statement}
\emph{Model-categories problem 8 (Mark Hovey).} Define a notion of \(2\)-model category and show that the intended \(2\)-category of model categories is an example. The proposed weak equivalences are at two levels: Quillen equivalences between model categories, and natural weak equivalences between Quillen functors. The source's belief that an ordinary model-category structure is insufficient is motivation, not a proved nonexistence theorem.

A precise formulation must choose objects and size universe, \(1\)-morphisms (for example, a specified handedness of Quillen functors or Quillen adjunctions), \(2\)-morphisms, and the required axioms. A Quillen equivalence is an adjunction inducing the intended equivalence of derived homotopy theories. A natural weak equivalence must be defined with its evaluation domain: objectwise equivalences on all objects and equivalences only on cofibrant or fibrant objects are different conventions. The source does not fix these choices or propose explicit \(2\)-model axioms.

The intended homotopy \(2\)-category should arise, in a higher universe, by localizing at the Quillen equivalences and the selected natural weak equivalences. The source says this localization at least makes sense and points to higher-category localization work. A completed definition must specify its localization meaning and universal property, rather than merely rename that formal inversion as a verified \(2\)-model structure.

The supplied summary describes infinity-categorical frameworks as modern successors and says a literal strict \(2\)-model realization was not identified, but its final sentence is truncated. Its identification with the entire infinity-category of infinity-categories needs object-class, universe and presentation hypotheses, and an \((\infty,1)\) model does not automatically retain all noninvertible natural transformations intended as \(2\)-cells. No equivalence of these different frameworks is asserted without those distinctions.
\subsection{Short English statement}
Formulate higher model-category axioms that treat Quillen equivalences and natural weak equivalences at separate categorical levels, and verify them for a specified \(2\)-category of model categories. Localization and modern infinity-categorical replacements do not by themselves fix the requested axioms or \(2\)-morphism structure.
\subsection{Sources}
\begin{itemize}
\item Local UnsolvedMath registry, permanent ID 11100054, source code catalogue entry; proposer Mark Hovey, ``Model categories,'' item 8.
\item \emph{Hovey's Algebraic Topology Problem List}, surviving mirror: \url{https://www-users.cse.umn.edu/~tlawson/hovey/}. Recorded access 2026-07-29; extraction html-ordered-list. The mirror contains 65 top-level items across nine category pages.
\item The statement names Andre Hirschowitz and Carlos Simpson, ``descent pour les n-champs,'' with the locator ``on xxx'' but no actual identifier. The summary names To\"en--Vezzosi, Lurie, and quasi-category approaches without exact references or theorem hypotheses.
\item Supplied status partially\_solved and PARTIAL-PROGRESS. The summary calls infinity-categories a partial modern realization rather than a literal strict \(2\)-model answer; its final discussion is truncated. Current primary-literature status and redistribution terms remain NEEDS\_REVIEW.
\end{itemize}
\end{document}
